
%(BEGIN_QUESTION)
% Copyright 2011, Tony R. Kuphaldt, released under the Creative Commons Attribution License (v 1.0)
% This means you may do almost anything with this work of mine, so long as you give me proper credit

Les og lag en disposisjon av Case History \#68 (``Control Mayhem On A Mine'') fra Michael Browns samling av veiledninger om optimalisering av reguleringssløyfer. Forbered deg på å diskutere begrepene og eksemplene i denne teksten grundig med lærer og medstudenter, og svar på følgende spørsmål:

\begin{itemize}
\item{} Undersøk PV- og Output (PD)-trendene i figur 1, og forklar hvor du mener ``trinnene'' i begge signalene kommer fra.
\vskip 10pt
\item{} Undersøk reguleringsstrategien for malmknuseren i figur 2, og forklar metallurgenes begrunnelse for ``selector''-strategien mellom nivåregulatoren og effektregulatoren.
\vskip 10pt
\item{} Beskriv den spesielle ``interlock''-strategien for nivå som var programmert i denne PLS-en, og som førte til at matebåndet ble stoppet automatisk. Beskriv også hvilken effekt dette hadde på PID-blokken mens den ble overstyrt av interlock-logikken.
\vskip 10pt
\item{} Et annet problem i knuserens nivåsløyfe var at hastigheten på matebåndet ikke fulgte regulatorutgangen nøyaktig. Forklar hvordan dette kan sees i trenden i figur 5, og hvordan dette er analogt med en feil dimensjonert reguleringsventil.
\vskip 10pt
\item{} I ``Addendum''-delen av casehistorien beskriver Brown et annet PLS-programmeringsproblem der en nødvendig funksjon i PID-blokkene ikke var tilgjengelig. Identifiser hvilken funksjon som var blokkert, hvorfor den er viktig, og hvorfor noen PLS-programmerer i det hele tatt kunne mene at det var greit å gjøre dette.
\end{itemize}

\vskip 20pt \vbox{\hrule \hbox{\strut \vrule{} {\bf Forslag til sokratisk diskusjon} \vrule} \hrule}

\begin{itemize}
\item{} Hvordan kan vi se at reguleringsventilen i slurrysløyfen (figur 1, lukket sløyfe-test) er overdimensjonert?
\item{} Identifiser hvilke tidspunkter i trenden i figur 3 der reguleringssystemet går i ``nivåmodus'' og hvilke tidspunkter det går i ``effektmodus''.
\item{} Kommenter {\it installert karakteristikk} for sluttelementet (f.eks. hurtigåpnende, lineær eller likeprosent) i sløyfen med trend vist i figur 5. Finnes det noe ``stiction'' i dette sluttelementet?
\end{itemize}


\underbar{file i01597}
%(END_QUESTION)





%(BEGIN_ANSWER)


%(END_ANSWER)





%(BEGIN_NOTES)

``Trinnene'' i signalene stammer mest sannsynlig fra treg scan-tid, enten i PLS-en eller eventuelt i transmitteren.

\vskip 10pt

Reguleringsstrategi for knuser: reguleringen vekslet fra nivå til motoreffekt hver gang effekten passerte en terskel (motorene var underdimensjonert). Resultatet var at reguleringen fortsatte å veksle mellom nivå og effekt i en kontinuerlig syklus. Personellet på gruven hadde forsøkt å ``tune'' sløyfen uten hell i årevis.

\vskip 10pt

En interlock-rutine tvang matebåndet til å stoppe dersom nivået i knuseren ble for høyt, men nivåregulatoren sto fortsatt i automatisk modus og fortsatte derfor å integrere.

\vskip 10pt

Figur 5 viser et ``overdimensjonert'' matebånd: båndhastigheten varierte mellom 1\% og 100\% mens regulatorutgangen varierte mellom 18\% og 38\%. Dette tilsvarer en prosessforsterkning på 5, som er svært høy. Kilden til problemet var beslutningen (tatt av en av metallurgene ved gruven) om å overdimensjonere den hydrauliske reguleringsventilen som strupet hydraulikkvæske til hydraulikkmotoren på transportbåndet.

\vskip 10pt

I et sørafrikansk malingsanlegg med tyske (Siemens?) PLS-er kunne PID-blokkene ikke settes i manuell modus for tuning. Denne reguleringen var dessuten kritisk: avvik på 1 grad Celsius betydde at batchen måtte kasseres.







%INDEX% Reading assignment: Michael Brown Case History #68, "Control mayhem on a mine"

%(END_NOTES)
